% Options for packages loaded elsewhere
\PassOptionsToPackage{unicode}{hyperref}
\PassOptionsToPackage{hyphens}{url}
%
\documentclass[
  10pt,
  a4paper,
]{article}
\usepackage{amsmath,amssymb}
\usepackage{iftex}
\ifPDFTeX
  \usepackage[T1]{fontenc}
  \usepackage[utf8]{inputenc}
  \usepackage{textcomp} % provide euro and other symbols
\else % if luatex or xetex
  \usepackage{unicode-math} % this also loads fontspec
  \defaultfontfeatures{Scale=MatchLowercase}
  \defaultfontfeatures[\rmfamily]{Ligatures=TeX,Scale=1}
\fi
\usepackage{lmodern}
\ifPDFTeX\else
  % xetex/luatex font selection
  \setmainfont[]{DejaVu Serif}
  \setsansfont[]{DejaVu Sans}
\fi
% Use upquote if available, for straight quotes in verbatim environments
\IfFileExists{upquote.sty}{\usepackage{upquote}}{}
\IfFileExists{microtype.sty}{% use microtype if available
  \usepackage[]{microtype}
  \UseMicrotypeSet[protrusion]{basicmath} % disable protrusion for tt fonts
}{}
\makeatletter
\@ifundefined{KOMAClassName}{% if non-KOMA class
  \IfFileExists{parskip.sty}{%
    \usepackage{parskip}
  }{% else
    \setlength{\parindent}{0pt}
    \setlength{\parskip}{6pt plus 2pt minus 1pt}}
}{% if KOMA class
  \KOMAoptions{parskip=half}}
\makeatother
\usepackage{xcolor}
\usepackage[left=15mm,right=15mm,top=20mm,bottom=20mm]{geometry}
\usepackage{longtable,booktabs,array}
\usepackage{calc} % for calculating minipage widths
% Correct order of tables after \paragraph or \subparagraph
\usepackage{etoolbox}
\makeatletter
\patchcmd\longtable{\par}{\if@noskipsec\mbox{}\fi\par}{}{}
\makeatother
% Allow footnotes in longtable head/foot
\IfFileExists{footnotehyper.sty}{\usepackage{footnotehyper}}{\usepackage{footnote}}
\makesavenoteenv{longtable}
\usepackage{graphicx}
\makeatletter
\def\maxwidth{\ifdim\Gin@nat@width>\linewidth\linewidth\else\Gin@nat@width\fi}
\def\maxheight{\ifdim\Gin@nat@height>\textheight\textheight\else\Gin@nat@height\fi}
\makeatother
% Scale images if necessary, so that they will not overflow the page
% margins by default, and it is still possible to overwrite the defaults
% using explicit options in \includegraphics[width, height, ...]{}
\setkeys{Gin}{width=\maxwidth,height=\maxheight,keepaspectratio}
% Set default figure placement to htbp
\makeatletter
\def\fps@figure{htbp}
\makeatother
\setlength{\emergencystretch}{3em} % prevent overfull lines
\providecommand{\tightlist}{%
  \setlength{\itemsep}{0pt}\setlength{\parskip}{0pt}}
\setcounter{secnumdepth}{-\maxdimen} % remove section numbering
\usepackage{sectsty}
\allsectionsfont{\raggedright}
\usepackage{lineno}
\linenumbers
\usepackage{setspace}
\doublespacing
\usepackage{fancyhdr}
\pagestyle{fancy}\fancyhf{}
\fancyhead[L]{\small Optical clocks distinguish neighbouring semidiurnal tides}
\fancyfoot[C]{\thepage}\setlength{\headheight}{14pt}
\usepackage{float}\floatplacement{figure}{H}
\setlength{\emergencystretch}{2em}
\ifLuaTeX
  \usepackage{selnolig}  % disable illegal ligatures
\fi
\IfFileExists{bookmark.sty}{\usepackage{bookmark}}{\usepackage{hyperref}}
\IfFileExists{xurl.sty}{\usepackage{xurl}}{} % add URL line breaks if available
\urlstyle{same}
\hypersetup{
  hidelinks,
  pdfcreator={LaTeX via pandoc}}

\author{}
\date{}

\begin{document}

\hypertarget{optical-clocks-distinguish-neighbouring-semidiurnal-tides}{%
\section{Optical clocks distinguish neighbouring semidiurnal
tides}\label{optical-clocks-distinguish-neighbouring-semidiurnal-tides}}

{[}Author list, order and affiliation assignments to be confirmed by the
authors{]}

\textsuperscript{1} Hefei National Laboratory, University of Science and
Technology of China, Hefei, China

\textsuperscript{2} Shanghai Research Center for Quantum Sciences and
CAS Center for Excellence in Quantum Information and Quantum Physics,
University of Science and Technology of China, Shanghai, China

\textsuperscript{3} Hefei National Research Center for Physical Sciences
at the Microscale and School of Physical Sciences, University of Science
and Technology of China, Hefei, China

\textsuperscript{4} Innovation Academy for Precision Measurement Science
and Technology, Chinese Academy of Sciences, Wuhan, China

\textsuperscript{5} Jinan Institute of Quantum Technology and CAS Center
for Excellence in Quantum Information and Quantum Physics, University of
Science and Technology of China, Jinan, China

\textsuperscript{6} Time and Frequency Geodesy Center (TFGC), School of
Geodesy and Geomatics, Wuhan University, Wuhan, China

\hypertarget{abstract}{%
\subsection{Abstract}\label{abstract}}

Optical clocks can track changes in Earth's gravitational potential, but
tides with similar periods are difficult to distinguish in short
comparisons.\textsuperscript{1,2,3} Here we distinguish two semidiurnal
responses between coastal Shanghai and inland Wuhan. The principal lunar
tide takes about 25 minutes longer per cycle than the principal solar
tide. Our 546 hours of observations, spread over 353 days, sample the
tides at different relative phases. The measured frequency ratio shows
distinct spectral peaks near both principal frequencies. Relative to the
combined solid-Earth and ocean-loading prediction, fitted amplitude
scales are \(1.00\pm0.13\) for the lunar group and \(0.93\pm0.36\) for
the neighbouring group containing the principal solar tide (one
statistical standard error). The changing lunar--solar phase provides a
reference for comparing tidal signals across interruptions even when the
frequency offsets between observing runs are unknown.

\hypertarget{main-text}{%
\subsection{Main text}\label{main-text}}

Two tides can look similar during a short observation, and interruptions
can leave an unknown frequency offset between runs. The principal lunar
semidiurnal tide M₂ repeats every 12.42 h and the principal solar tide
S₂ every 12 h. Observations at different lunar--solar phases can
therefore distinguish their responses without joining the runs at a
common frequency level.

We use this phase change in a comparison of a strontium clock in
Shanghai and a ytterbium clock in Wuhan. The record contains 545.9 h of
accepted observations on 43 UTC dates spread over 353 days. A joint
model with an independent constant for every valid segment resolves two
semidiurnal groups, M₂+N₂ and S₂+K₂. The observed spectrum also has a
local peak near each principal frequency. K₂ contains both lunar and
solar forcing; the neighbouring group therefore contains more than the
purely solar S₂ tide. The group's internal phases and relative
amplitudes follow the astronomical prediction. The approximately daily
constituents K₁, O₁, P₁ and Q₁ are estimated together as a third group.

Optical clocks measure potential differences through their relative
rates.\textsuperscript{1,4} Clock comparisons have resolved laboratory
height-dependent redshifts and enabled remote geodetic
measurements.\textsuperscript{5,6,7} Stabilized optical-fibre networks
connect clocks without transporting them.\textsuperscript{8,9}
Time-dependent measurements must distinguish gravitational responses
from changes in the clocks and transfer system. Astronomical tides
supply frequencies and phases predicted independently of the clock
observations.\textsuperscript{2,10,11}

\hypertarget{observing-dates-and-the-coastalinland-baseline}{%
\subsubsection{Observing dates and the coastal--inland
baseline}\label{observing-dates-and-the-coastalinland-baseline}}

The two clocks are 692 km apart and connected through Hefei by a
1,350-km phase-stabilized fibre link (Fig. 1a).\textsuperscript{12}
Between September 2025 and August 2026 they supplied 3,299 ten-minute
means in 54 valid segments and eight observing periods (Fig. 1c).
Observations occur in eight calendar months, with actual UTC times and
all missing intervals retained. The June--August periods contribute
273.9 h and sample different lunar--solar phases; the earlier periods
add further phases and operating conditions (Fig. 1d). The placement of
observations in the lunar--solar cycle supplies the separation between
the nearby signals.

The earlier summer comparison reports a nominal white-frequency-noise
trend of \(5.4\times10^{-16}/\sqrt{\tau/\mathrm{s}}\), where \(\tau\) is
averaging time.\textsuperscript{12} This describes the linked-clock
system's technical performance, while the present amplitude errors are
estimated from the observed record. The individual Sr and Yb systems are
described in their clock evaluations.\textsuperscript{13,14} Some summer
observations overlap the absolute frequency-ratio measurement of Hou and
colleagues; here the measured quantity is the time-dependent tidal
response. A positive fractional ratio variation means Wuhan Yb runs
faster relative to Shanghai Sr.~In the geodetic potential convention,
\(y=-[\delta W_{\mathrm{Wuhan}}-\delta W_{\mathrm{Shanghai}}]/c^2\).

Tides alter the direct gravitational potential, deform the Earth and
move the clocks. Ocean tides add the attraction of redistributed water,
the Earth's induced potential and displacement of the clock. The
prediction combines an IERS-based solid tide with FES2022 ocean
elevations convolved with a kernel containing these three loading
contributions (Methods).\textsuperscript{15,16,17,18} Shanghai lies
beside the East China Sea and Wuhan inland, making their differential
potential sensitive to ocean loading.

For M₂, the predicted solid contribution is \(3.43\times10^{-18}\) and
the loading contribution \(1.42\times10^{-18}\) in fractional clock
rate. Their phases differ by 19.8°, giving a combined amplitude of
\(4.78\times10^{-18}\) (Fig. 1b). After removing segment means,
loading's weighted root-mean-square variation is 33.1\% of the combined
prediction on the sampled record. This model contribution is distinct
from an observed loading coefficient. A separate filtered analysis
allowing loading its own coefficient reported \(0.89\pm0.55\); that
coefficient remains unresolved (Methods).

The zero-point choice follows the available experimental records. Of the
segment boundaries within the eight periods, 46 have no documented
calibrated zero transfer. We therefore allow each of the 54 segments its
own constant, retaining the changing tidal signal within each run.
Physical responses are estimated jointly from the original ten-minute
means with fixed noise weights. No empirical daily sinusoid is removed
before this fit, and no fitted values fill the gaps.

\hypertarget{two-semidiurnal-responses-and-two-spectral-peaks}{%
\subsubsection{Two semidiurnal responses and two spectral
peaks}\label{two-semidiurnal-responses-and-two-spectral-peaks}}

The joint fit gives \(A_L=1.00\pm0.13\) for M₂+N₂ and
\(A_S=0.93\pm0.36\) for S₂+K₂ (Fig. 2b). One denotes the combined
solid-Earth and ocean-loading prediction. Errors are one statistical
standard error, taken as the larger of residual-scaled least squares and
UTC-day deletion; components and conditional coverage tests are given in
Methods. The daily group is estimated in the same model as
\(A_D=0.58\pm0.54\). Its larger uncertainty is propagated through the
joint fit, which still constrains both semidiurnal groups.

The same measured frequency-ratio record shows local spectral peaks at
22.361 and 23.147 µHz, close to the prescribed M₂ and S₂ frequencies of
22.364 and 23.148 µHz (Fig. 2a). The scan uses all 3,299 observed means
at their actual UTC times, with 54 segment constants and equal weights.
The peak separation exceeds the inverse calendar span, 0.0328 µHz. The
left axis gives density in gravitational-potential units and the right
the equivalent fractional-frequency density. A fractional variation of
\(10^{-18}\) corresponds to about \(0.09\,\mathrm{m^2/s^2}\) in
potential difference.

The observation window also produces sidelobes. Sampling a single tidal
line only at the measured times spreads power towards neighbouring
frequencies. Passing individual M₂ and S₂ predictions through this same
sampling and scan gives 8.03\% M₂-to-S₂ cross-power and 8.15\% in the
reverse direction. These are responses of the sampling window, not
random noise generated by missing observations. A conditional scan
retains a local feature near each principal frequency after projecting
out the sine and cosine at the other (Extended Data Fig. 2).

These checks support assigning the two features to the principal lunar
and solar semidiurnal neighbourhoods. The grouped amplitudes come from
the simultaneous physical templates, with prescribed internal
relationships and Allan-based noise weights; they do not treat N₂ or K₂
as individually resolved spectral lines. The spectrum, including its
curve after subtracting the fixed unit prediction, displays the same
observations and supplies no independent detection probability.

The phase of the recurring signal is also compatible with the
prediction. Adding a lunar component one quarter-cycle from the
predicted phase gives in-phase and quadrature responses of
\(1.00\pm0.13\) and \(-0.10\pm0.20\) (Fig. 2c). Figure 3 orders the full
record into twelve M₂ phase bins. The observations after subtracting the
fixed minor remainder, and the unit lunar prediction, undergo identical
projection against the segment constants and other fitted groups.
Propagated errors retain nuisance-parameter uncertainty and covariance
between phase means. All individual points and each bin's support remain
visible, so the folded curve can be traced to the full record.

\hypertarget{support-across-observing-periods-and-instrument-controls}{%
\subsubsection{Support across observing periods and instrument
controls}\label{support-across-observing-periods-and-instrument-controls}}

The lunar response is supported on both sides of summer. With the same
segment model it is \(1.21\pm0.23\) before summer and \(0.77\pm0.18\) in
summer. Deleting one observing period at a time gives point estimates
from 0.885 to 1.093 (Fig. 4a). These overlapping fits show how each
period contributes to the full-record estimate.

As a same-data check, sharing one reference level within each of the
eight periods gives a lunar coefficient of \(0.987\pm0.123\) and a
neighbouring coefficient of \(1.129\pm0.288\). Only the intercept
grouping changes in this comparison; observations, physical templates
and weights are identical. Residual structure and deletion results are
retained in Methods and Extended Data Fig. 1. The segment model is
selected from the zero-transfer evidence, rather than the proximity of a
coefficient to one.

A simultaneous fibre round-trip record tests the monitored transmission
channel on 246.6 h of summer observations, from 29 June to 25 August
2026. Clock and monitor means contain matching valid members. Adding the
monitor to the otherwise identical fit changes the lunar coefficient by
\(+0.0061\pm0.0063\) and the neighbouring coefficient by
\(+0.0059\pm0.0145\) (Fig. 4b). Its coverage is confined to that overlap
and transmission route; other local transfer branches and clock effects
require their own observations.

Temperature changes and operating schedules motivate checks at 24 h. A
separate filtered analysis reports \(0.98\pm0.13\) with two-site
temperature regressors, whose sampling and supporting measurements are
described in Methods. Adding a shared 24-h sine--cosine pair on the
present observations gives \(1.04\pm0.13\). The weakly constrained
physical daily group is not a bound on laboratory disturbance. Freely
phased 12-h terms, which compete with the astronomical semidiurnal
groups within short runs, are retained only as sensitivity tests in
Methods.

The astronomical phase continues across an interruption even when the
measured frequency zero is not transferred. This phase reference allows
a sparse clock comparison to distinguish nearby tidal responses across a
353-day span. Contemporaneous European clock-network observations by
Stahl and colleagues track time-dependent potential on different
baselines and with different tidal mixtures.\textsuperscript{3} They
complement the grouped semidiurnal measurement on this coastal--inland
link. Longer coastal records at lower noise, with phase coverage that
distinguishes the solid and loading contributions and calibrated local
transfer paths, could allow ocean loading to be tested separately.

\clearpage

\singlespacing

\hypertarget{figure-legends}{%
\subsection{Figure legends}\label{figure-legends}}

\includegraphics{figures/figure1.pdf}

\textbf{Fig. 1 \textbar{} Observing geometry and lunar--solar phase
coverage.} \textbf{a,} Shanghai Sr and Wuhan Yb, 692 km apart, with the
1,350-km fibre connection through Hefei drawn schematically. \textbf{b,}
Predicted differential M₂ solid and ocean-loading phasors, with the
solid component defining phase zero; amplitudes are fractional clock
rates in units of \(10^{-18}\). \textbf{c,} Actual UTC sampling on 43
dates over 353 days, in eight observing periods, with accepted nominal
member-hours (545.9 h in total). \textbf{d,} The S₂−M₂ astronomical
phase at those same samples, including nodal corrections. Summer periods
P6--P8 are orange. Geography and phasors describe the prediction; the
lower panels describe the observation window.

\clearpage

\includegraphics{figures/figure2.pdf}

\textbf{Fig. 2 \textbar{} Two semidiurnal spectral peaks and their
grouped responses.} \textbf{a,} Equal-weight scan of all 3,299 observed
means at their true UTC times, with 54 segment constants. Peaks occur at
22.361 and 23.147 µHz, near M₂ and S₂. Blue is the observed scan; orange
follows subtraction of the fixed unit total prediction. Gaps are
retained. The two axes express the same nominal one-sided projection
density per µHz, related by \(S_W=c^4S_y\). The actual sampling yields
about 8\% bidirectional single-line M₂/S₂ cross-power. \textbf{b,} Joint
physical-group amplitudes with one standard error, defined as the larger
of scaled least-squares and UTC-day-jackknife errors. These estimates
use the physical templates and Allan-based noise weights. \textbf{c,}
Lunar in-phase and quadrature responses; the cross marks the unit
prediction. Contours have scaled-least-squares covariance radii one and
two, not separately calibrated joint confidence levels. The spectral
peaks identify principal-frequency neighbourhoods, while \textbf{b}
estimates prescribed constituent groups. All panels use the same
observation rows and are not independent significance tests.

\clearpage

\includegraphics{figures/figure3.pdf}

\textbf{Fig. 3 \textbar{} The full record ordered by lunar phase.}
\textbf{a,} Twelve 30° M₂ phase means of all 3,299 observations, after
subtracting the fixed minor remainder and projecting against 54 segment
constants and the neighbouring semidiurnal and daily groups. Identical
weighted projection and averaging are applied to the unit lunar
prediction; the fitted curve uses the unbinned primary coefficient.
Error bars are one working standard error with the nuisance projection
propagated; phase means are correlated. \textbf{b,} Every projected
observation, with the full observed vertical range. \textbf{c,} Number
of rows in each bin. Data, weights and templates are the primary segment
design throughout. This is a display of the main observations, not an
independent experiment.

\clearpage

\includegraphics{figures/figure4.pdf}

\textbf{Fig. 4 \textbar{} Temporal support and a simultaneous
transmission control.} \textbf{a,} Lunar estimates from the full record,
the pre-summer and summer subsets, and deletion of each observing
period. Each fit uses segment constants and the primary noise weights,
with errors recalculated on its retained rows. The shaded band is the
full-record one-standard-error range. These fits share observations.
\textbf{b,} Changes caused by including the simultaneous round-trip
monitor, with paired UTC-day-jackknife standard errors. Both fits use
identical clock/monitor members: 1,482 means representing 246.6 h in
P6--P8, between 29 June and 25 August 2026. This control covers the
monitored transmission channel on its measured overlap; the principal
tide fit retains the entire 545.9-h clock record.

\clearpage

\doublespacing

\hypertarget{methods}{%
\subsection{Methods}\label{methods}}

\hypertarget{observations-frequency-sign-and-averaging}{%
\subsubsection{Observations, frequency sign and
averaging}\label{observations-frequency-sign-and-averaging}}

The Shanghai and Wuhan coordinates used by the frozen prediction are
31.130° N, 121.570° E and 30.529° N, 114.357° E. The measured ratio is
Wuhan Yb/Shanghai Sr.~With positive geodetic potential \(W\), its tidal
fractional variation is
\(-[\delta W_{\mathrm{Wuhan}}-\delta W_{\mathrm{Shanghai}}]/c^2\). The
period-relative beat conversion is
\(-d_p(f_{\mathrm{beat}}-f_{{\mathrm{ref}},p})/F_{1550}\), where
\(F_{1550}=1.933992\times10^{14}\) Hz and \(d_p=2\) in 2025 and 1 in
2026. Thus a 2026 beat deviation of 1 Hz has fractional magnitude
\(5.17\times10^{-15}\). The 2025 division factor is applied before
forming the small clock-rate variation.

The supplied valid-interval table defines 54 segments in 29 operating
groups and eight periods. Counter members are read across midnight
before binning. Beijing-time labels are converted to UTC. Summer group
corrections follow the decimal-arithmetic frequency-chain calculation;
early absolute correction entries remain incomplete. Their constant part
is represented by segment intercepts, while any time variation within a
segment remains a possible systematic contribution. Fitted intercepts
are never applied to join the runs for long-term stability estimation.

The membership audit contains 1,966,908 nominal one-second rows before
the minimum-bin-support cut. Three boundary rows reconcile its
difference from the printed interval total. Means are formed in nominal
600-s bins separately within each segment and require at least 180
retained members. The 3,299 accepted means contain 1,965,327 members, or
545.924 h. There are 3,291 unique bin centres: eight centres contain
separate partial means from different segments, with disjoint members.
The 353.39-day span contains observations on 43 UTC dates. The inherited
valid mask is retained; this work neither applies an additional
triangular filter nor selects points using tidal residuals. Exact
raw-file memberships and hashes are retained in the controlled archive.

\hypertarget{what-the-constituent-names-mean}{%
\subsubsection{What the constituent names
mean}\label{what-the-constituent-names-mean}}

Semidiurnal tides produce approximately two cycles per day. M₂ is the
principal lunar constituent; N₂ describes an additional lunar variation
associated with orbital ellipticity. S₂ is the principal solar
semidiurnal constituent, whereas K₂ contains lunar and solar forcing.
The names identify astronomical frequencies, not separate instruments or
independent experiments. The diurnal group contains K₁ (lunisolar), O₁
(principal lunar), P₁ (principal solar) and Q₁ (lunar elliptic). Their
periods are listed in Extended Data Table 3.\textsuperscript{19}
Internal relative amplitudes, phases and nodal modulation are fixed
within each physical group. A group response therefore tests that
complete specified pattern.

\hypertarget{combined-solid-and-loading-potential}{%
\subsubsection{Combined solid and loading
potential}\label{combined-solid-and-loading-potential}}

The physical target includes the direct astronomical potential, the
Earth's induced potential and the effect of clock displacement. The
selected frozen solid sequence implements these through the IERS radial
displacement \(u\) and a fixed effective-potential conversion,

\[y_{\mathrm{solid}}=-\frac{1+k_2-h_2}{h_2}\frac{g(u_{\mathrm{Wuhan}}-u_{\mathrm{Shanghai}})}{c^2},\]

with \(h_2=0.6078\), \(k_2=0.3010\) and \(g=9.794\,\mathrm{m/s^2}\).
This is a displacement-based approximation to the combined solid
potential, not the alternative frequency-resolved direct-potential
implementation. The same conversion also acts on degree-three and
frequency-dependent displacement contributions in the supplied sequence.
A previous implementation comparison found an approximately 18\% K₁
amplitude deficit relative to a potential-weighted reference. The daily
response is accordingly not used to infer a Love number. Uncertainty in
the solid model remains separate from the quoted statistical errors.

For loading, FES2022 sea-surface elevations \(\zeta\) are convolved with
the effective-potential kernel

\[\delta W_{\mathrm{load}}=\int \rho_w\zeta(\mathbf{r}',t)G_W(\psi)\,\mathrm{d}A,\qquad
G_W(\psi)=\frac{ga}{M_E}\sum_n(1+k'_n-h'_n)P_n(\cos\psi).\]

The three terms represent direct water attraction, induced potential and
clock displacement. PREM load Love numbers are retained through degree
696 with the archived high-degree continuation. Eighteen loading
constituents are synthesized with their astronomical nodal corrections.
The complex convention is \(f\,\mathrm{Re}[C\exp(-i\phi)]\); reversing
the exponential sign requires conjugating \(C\). The present analysis
reads the frozen constituent and time-series tables rather than
substituting an older loading sequence.

The eight named constituents are grouped as described above. Minor terms
remain once as a fixed remainder \(r_0\). On the actual weighted
segment-centred sampling, solid, loading and total root-mean-square
variations are respectively 2.928, 1.328 and \(4.011\times10^{-18}\).
Their ratio 0.331 is an amplitude-norm ratio, not an additive fraction
of total power. The separately reported \(0.89\pm0.55\) loading
coefficient belongs to the alternative directly modelled, filtered
reduction described below; it has not been recalculated as a separate
coefficient of the present estimator.

\hypertarget{joint-estimator-and-zero-point-choice}{%
\subsubsection{Joint estimator and zero-point
choice}\label{joint-estimator-and-zero-point-choice}}

We fit

\[y_i=c_{s(i)}+A_Ls_{L,i}+A_Ss_{S,i}+A_Ds_{D,i}+r_{0,i}+\epsilon_i,\]

with 54 segment intercepts and three physical response coefficients.
Weighted least squares uses the frozen Allan-based uncertainties
\(\sigma_i\), with \(W_{ii}=1/\sigma_i^2\). A period noise coefficient
\(\alpha_p\) is estimated from within-segment overlapping Allan
deviations at 256, 512, 1,024, 2,048 and 4,096 s, after removing the
fixed unit tide. No Allan window crosses a segment boundary. The bin
uncertainty is \(\sigma_i=\alpha_{p(i)}/\sqrt{n_i}\), where \(n_i\) is
its retained member count. This is a working noise model, not stability
calculated after joining fit-adjusted levels. Independent QR replay
reproduces the eight-period design before changing only its intercept
grouping. Sampling, corrections, templates and weights are held fixed.

The inspected interval and correction records contain no documented
calibrated zero transfer for the 46 adjacent segment boundaries that
share an eight-period intercept. They give valid times, but not a signed
boundary transfer with uncertainty or a lock/relock history establishing
unchanged reference levels. This motivates the segment design. The
eight-period model's weighted residual sum of squares is 3,781.58; the
segment model's is 3,713.64, with 46 additional intercept parameters.
Their unweighted residual root-mean-square values are 24.114 and
\(23.934\times10^{-18}\). Reduced residual sums of squares are 1.1501
and 1.1455. The numerical improvement alone does not establish a
physical jump at each boundary.

With period constants, weighted segment-mean residuals range from
\(-6.75\) to \(12.00\times10^{-18}\). Adjacent standardized residual
correlations, restricted to the same segment and separations of at most
900 s, are −0.127 and −0.148 for the period and segment designs.
Deleting a period gives lunar point-estimate ranges 0.877--1.052 and
0.885--1.093, respectively (Extended Data Fig. 1).

\hypertarget{statistical-errors-and-conditional-coverage}{%
\subsubsection{Statistical errors and conditional
coverage}\label{statistical-errors-and-conditional-coverage}}

For each group the reported standard error is the larger of
residual-scaled least squares and leave-one-UTC-day-out jackknife
errors. Deletions remove a calendar day, refit all remaining
coefficients and drop empty intercept columns. The jackknife variance is
\((G-1)G^{-1}\sum_g(\hat A_{(-g)}-\overline A)^2\). Period-deletion
errors are recorded separately. Extended Data Table 1 gives all
components and the same-data intercept comparisons. Observations are
repeated measurements of one clock pair, not thousands of independent
experiments.

New conditional simulations apply both designs to the same 3,299-row
sampling and heteroscedastic uncertainties. Each of four scenarios has
1,000 realizations: independent Gaussian noise; equal variance
contributions from Gaussian white noise and a 6-h Ornstein--Uhlenbeck
process within segments; the corresponding 24-h process within periods;
and standardized Student-\(t_5\) noise. The overall scale is the
observed segment-model residual scale; all true physical response
coefficients are one. Each realization recomputes both components of the
standard-error rule. Seed 2026110617 fixes the simulation. Coverage of
\(\hat A\pm1.96u\) is 95.1--97.4\% for the segment-model lunar
coefficient and 95.4--96.8\% for its neighbouring coefficient; Monte
Carlo standard errors are about 0.4--0.7 percentage points (Extended
Data Table 2). These are tests under declared noise laws, not a
calibration against every possible disturbance.

The phase-relaxed model adds the lunar quadrature while retaining the
prescribed M₂/N₂ relationship. Covariance contours use scaled least
squares. Subset and deletion fits refit the same model on their retained
rows. No shared-data estimates are averaged, and no half-daily
coefficient is converted into a daily upper limit.

\hypertarget{phase-means-spectral-density-and-sampling-sidelobes}{%
\subsubsection{Phase means, spectral density and sampling
sidelobes}\label{phase-means-spectral-density-and-sampling-sidelobes}}

For the phase plot, \(N\) contains segment constants and the two
non-lunar groups. We project both \(y-r_0\) and the unit lunar
prediction with \(M_N=I-N(N^TWN)^{-1}N^TW\), and average in twelve fixed
M₂ phase bins. If \(B\) is weighted bin averaging, the mean covariance
is \(BM_NC_yM_N^TB^T\), where \(C_y\) is the scaled working observation
covariance. It retains correlations between means and propagates
nuisance estimation. The full unbinned response remains primary.

Figure 2a preserves the independently specified equal-weight,
54-intercept scan of the unchanged observed means, rather than replacing
it with the amplitude fit's weights. Every trial sine--cosine pair and
the data are projected against the same constants. For a reduction in
squared residuals \(\Delta\mathrm{RSS}(f)\) with dimensionless \(y\),
the nominal one-sided projection density is
\(S_y(f)=600\,\mathrm{s}\,\Delta\mathrm{RSS}(f)\) per Hz, or \(10^{-6}\)
times that value per µHz. Potential density is \(S_W=c^4S_y\); amplitude
density is \(c^2\sqrt{S_y}\). Uniform-grid checks reproduce the
conventional one-sided Fourier density. For interrupted data this
normalization is a stated projection density, not a Parseval-calibrated
variance over a filled year.\textsuperscript{20}

The frequency step is 0.002729 µHz; the inverse calendar span is
0.032752 µHz. Fine plotting steps do not improve physical resolution.
M₂-only and S₂-only astronomical injections use the same actual times,
nodal factors and equal-weight segment projection. Power at the opposite
principal frequency, divided by power at the injected frequency, gives
0.08027 and 0.08149. Conditional scans project out a freely phased
sinusoid at the other frequency before scanning the target
neighbourhood. Amplitude interference means that such corrections act on
sampled waveforms or complex coefficients, not by subtracting scalar
power spectra. Gaps alone do not predict the realization of measurement
noise. Residual density can describe remaining measurement fluctuations,
but includes possible model mismatch and cannot be called a pure
instrumental noise floor.

\hypertarget{instruments-environmental-checks-and-the-alternative-reduction}{%
\subsubsection{Instruments, environmental checks and the alternative
reduction}\label{instruments-environmental-checks-and-the-alternative-reduction}}

The apparatus performance coefficient quoted in the main text describes
the gap-concatenated summer record in the earlier frequency-ratio
manuscript.\textsuperscript{12} The present noise coefficients are
estimated separately using only within-segment Allan windows, as defined
above.

The matched monitor calculation uses 1,482 rows with 887,882 nominal
members, or 246.634 h, in P6--P8. Both baseline and augmented fits have
identical members, segment labels and weights. Paired day-deletion
replicates propagate uncertainty in their difference. The lunar,
neighbouring and daily coefficient changes are respectively
\(0.00606\pm0.00627\), \(0.00594\pm0.01453\) and \(0.00791\pm0.01388\).
The archive does not provide that paired control for the pre-summer
periods or every laboratory transfer branch.

The alternative v68 reduction reports \(A_L=0.97\pm0.14\) from
approximately 544 h in 54 segments, with a 1,200-s triangular filter
sampled every 600 s, equal weights, a direct-potential solid model and a
continuous-run circular-residual-shift error. It reports \(0.98\pm0.13\)
with two-site ERA5 temperature regressors and a 474-h round-trip overlap
changing lunar and whole-template coefficients by less than 0.005. These
are author-reported manuscript results, not the present replay. The
stated 98-h Sr chamber record is a distinct thermometer measurement;
ERA5 2-m air temperatures are environmental proxies. Exact membership
for the 474-h control and temperature-regression support must be
supplied for an independent replay. These checks are not combined with
the current uncertainty.

On the selected data and model, a common 24-h sine--cosine pair gives
\(A_L=1.036\pm0.133\); separate 24-h pairs in each period give
\(1.003\pm0.136\). These are mechanism-motivated daily sensitivity
designs, not measurements of an independently observed thermal
amplitude. Free 12-h pairs in each period give \(0.751\pm0.242\) and a
very poorly constrained neighbouring coefficient; the design competes
with astronomical semidiurnal variation within short periods. It is
retained as a sensitivity boundary, with no assumed laboratory
disturbance assigned to it. A sequential residual diagnostic must
project its test basis through the same physical-fit operator and
propagate uncertainty.

\singlespacing

\hypertarget{data-availability}{%
\subsection{Data availability}\label{data-availability}}

The author review package contains the 3,299 derived means, fixed
templates, noise weights, figure source tables and immutable input
hashes. The original counters and membership ledger are held in the
controlled experimental archive. {[}Authors to supply an
editor-accessible location and confirm the permanent data-access
statement and repository identifier before submission.{]}

\singlespacing

\hypertarget{code-availability}{%
\subsection{Code availability}\label{code-availability}}

The package contains the independent QR replay, same-data design
comparison, conditional simulations, spectral calculations, covariance
propagation and figure scripts, with software versions and a
reproducible entry point. {[}Authors to confirm the code licence and
permanent repository identifier before submission.{]}

\singlespacing

\hypertarget{author-statements}{%
\subsection{Author statements}\label{author-statements}}

\textbf{Author contributions:} {[}To be completed and approved by all
authors.{]}

\textbf{Acknowledgements and funding:} {[}To be completed by the
authors.{]}

\textbf{Competing interests:} {[}Declaration to be supplied by the
authors.{]}

\textbf{Correspondence:} {[}Corresponding author names and email
addresses to be supplied.{]}

\textbf{AI-assisted preparation:} OpenAI Codex assisted English
drafting, numerical replay and figure preparation. {[}Authors to confirm
their verification of the analyses, citations and final text, and
responsibility for the submitted work.{]}

\clearpage

\singlespacing

\hypertarget{references}{%
\subsection{References}\label{references}}

\begin{enumerate}
\def\labelenumi{\arabic{enumi}.}
\tightlist
\item
  Mehlstäubler, T. E., Grosche, G., Lisdat, C., Schmidt, P. O. \&
  Denker, H. Atomic clocks for geodesy. \emph{Rep.~Prog. Phys.}
  \textbf{81}, 064401 (2018).
  \url{https://doi.org/10.1088/1361-6633/aab409}
\item
  Voigt, C., Denker, H. \& Timmen, L. Time-variable gravity potential
  components for optical clock comparisons and the definition of
  international time scales. \emph{Metrologia} \textbf{53}, 1365--1383
  (2016). \url{https://doi.org/10.1088/0026-1394/53/6/1365}
\item
  Stahl, K. et al.~Optical clock network tracks tidal variations of the
  gravity potential. \emph{arXiv}, 2610.39635v1 (2026).
  \url{https://arxiv.org/abs/2610.39635v1}
\item
  Ludlow, A. D., Boyd, M. M., Ye, J., Peik, E. \& Schmidt, P. O. Optical
  atomic clocks. \emph{Rev.~Mod. Phys.} \textbf{87}, 637--701 (2015).
  \url{https://doi.org/10.1103/revmodphys.87.637}
\item
  Chou, C. W., Hume, D. B., Rosenband, T. \& Wineland, D. J. Optical
  Clocks and Relativity. \emph{Science} \textbf{329}, 1630--1633 (2010).
  \url{https://doi.org/10.1126/science.1192720}
\item
  McGrew, W. F. et al.~Atomic clock performance enabling geodesy below
  the centimetre level. \emph{Nature} \textbf{564}, 87--90 (2018).
  \url{https://doi.org/10.1038/s41586-018-0738-2}
\item
  Grotti, J. et al.~Geodesy and metrology with a transportable optical
  clock. \emph{Nat. Phys.} \textbf{14}, 437--441 (2018).
  \url{https://doi.org/10.1038/s41567-017-0042-3}
\item
  Lisdat, C. et al.~A clock network for geodesy and fundamental science.
  \emph{Nat. Commun.} \textbf{7}, 12443 (2016).
  \url{https://doi.org/10.1038/ncomms12443}
\item
  Riehle, F. Optical clock networks. \emph{Nat. Photon.} \textbf{11},
  25--31 (2017). \url{https://doi.org/10.1038/nphoton.2016.235}
\item
  Denker, H. et al.~Geodetic methods to determine the relativistic
  redshift at the level of \(10^{-18}\) in the context of international
  timescales: a review and practical results. \emph{J. Geod.}
  \textbf{92}, 487--516 (2018).
  \url{https://doi.org/10.1007/s00190-017-1075-1}
\item
  Schröder, S., Stellmer, S. \& Kusche, J. Potential and scientific
  requirements of optical clock networks for validating
  satellite-derived time-variable gravity data. \emph{Geophys. J. Int.}
  \textbf{226}, 764--779 (2021).
  \url{https://academic.oup.com/gji/article/226/2/764/6212226}
\item
  Hou, L. et al.~Yb/Sr optical clock comparison at \(2.2\times10^{-18}\)
  uncertainty over a 1350-km fiber link. \emph{arXiv}, 2610.34250v1
  (2026). \url{https://arxiv.org/abs/2610.34250v1}
\item
  Jia, Z.-P. et al.~Improved systematic evaluation of a strontium
  optical clock with uncertainty below \(1\times10^{-18}\).
  \emph{Metrologia} \textbf{63}, 025002 (2026).
  \url{https://doi.org/10.1088/1681-7575/ae449e}
\item
  Zhu, Q. et al.~Ytterbium lattice clock with systematic uncertainty of
  \(1.3\times 10^{-18}\) and instability at the \(10^{-19}\) level.
  \emph{arXiv}, 2606.10514v2 (2026).
  \url{https://arxiv.org/abs/2606.10514v2}
\item
  Petit, G. \& Luzum, B. (eds) IERS Conventions (2010). \emph{IERS
  Technical Note} \textbf{36} (2010).
  \url{https://iers-conventions.obspm.fr/content/tn36.pdf}
\item
  CNES. FES2022 (Finite Element Solution) Tidal model (Version 2024).
  CNES \url{https://doi.org/10.24400/527896/A01-2024.004} (2024).
\item
  Farrell, W. E. Deformation of the Earth by surface loads.
  \emph{Rev.~Geophys.} \textbf{10}, 761--797 (1972).
  \url{https://doi.org/10.1029/RG010i003p00761}
\item
  Han, D. \& Wahr, J. The viscoelastic relaxation of a realistically
  stratified earth, and a further analysis of postglacial rebound.
  \emph{Geophys. J. Int.} \textbf{120}, 287--311 (1995).
  \url{https://doi.org/10.1111/j.1365-246x.1995.tb01819.x}
\item
  NOAA. Harmonic constituents. \emph{NOAA Tides and Currents}.
  \url{https://tidesandcurrents.noaa.gov/harcon.html?id=8577330}
  (accessed 6 November 2026).
\item
  VanderPlas, J. T. Understanding the Lomb--Scargle periodogram.
  \emph{Astrophys. J. Suppl. Ser.} \textbf{236}, 16 (2018).
  https://doi.org/10.3847/1538-4365/aab766
\end{enumerate}

\clearpage

\singlespacing\sffamily\fontsize{7}{9}\selectfont

\hypertarget{extended-data}{%
\subsection{Extended Data}\label{extended-data}}

\includegraphics{figures/ed_figure1.pdf}

\textbf{Extended Data Fig. 1 \textbar{} Same-data comparison of
reference-level designs.} \textbf{a,} Weighted mean residual of each
segment after the eight-period joint fit. The additional segment
intercepts make these means zero by construction, so a zero line would
not demonstrate calibration. \textbf{b,} Lunar point estimates after
deleting each observing period under eight-period and 54-segment
designs. Both use the same frozen observations, physical templates,
corrections and noise weights. These diagnostics do not replace
lock/relock or zero-transfer records.

\includegraphics{figures/ed_figure2.pdf}

\textbf{Extended Data Fig. 2 \textbar{} Conditional semidiurnal scans.}
Equal-weight scans of the same 3,299 observations with 54 segment
intercepts. \textbf{a,} M₂ neighbourhood before and after projection of
a sine--cosine pair at S₂. \textbf{b,} S₂ neighbourhood before and after
projection at M₂. Dotted lines show the astronomical target frequencies.
These are descriptive projections on shared observations, not additional
detection probabilities.

\clearpage

\textbf{Extended Data Table 1 \textbar{} Estimator and uncertainty
comparison.} All directly replayed designs use the same 3,299 means,
physical templates and Allan-based weights. Errors are standard errors.
The alternative filtered estimate is separately labelled
author-reported.

\begin{longtable}[]{@{}
  >{\raggedright\arraybackslash}p{(\columnwidth - 8\tabcolsep) * \real{0.2706}}
  >{\raggedright\arraybackslash}p{(\columnwidth - 8\tabcolsep) * \real{0.1765}}
  >{\raggedright\arraybackslash}p{(\columnwidth - 8\tabcolsep) * \real{0.1765}}
  >{\raggedright\arraybackslash}p{(\columnwidth - 8\tabcolsep) * \real{0.1765}}
  >{\raggedright\arraybackslash}p{(\columnwidth - 8\tabcolsep) * \real{0.2000}}@{}}
\toprule\noalign{}
\begin{minipage}[b]{\linewidth}\raggedright
Design
\end{minipage} & \begin{minipage}[b]{\linewidth}\raggedright
Lunar
\end{minipage} & \begin{minipage}[b]{\linewidth}\raggedright
Near\_12h
\end{minipage} & \begin{minipage}[b]{\linewidth}\raggedright
Diurnal
\end{minipage} & \begin{minipage}[b]{\linewidth}\raggedright
Status
\end{minipage} \\
\midrule\noalign{}
\endhead
\bottomrule\noalign{}
\endlastfoot
S54 & 1.000 ± 0.132 & 0.925 ± 0.365 & 0.579 ± 0.538 & Replayed \\
G29 & 0.984 ± 0.124 & 1.093 ± 0.298 & 0.759 ± 0.420 & Replayed \\
P8 & 0.987 ± 0.123 & 1.129 ± 0.288 & 0.678 ± 0.359 & Replayed \\
Filtered, 54 segments & 0.97 ± 0.14 & 0.90 ± 0.29 & 0.35 ± 0.51 &
Author-reported \\
\end{longtable}

\begin{longtable}[]{@{}
  >{\raggedright\arraybackslash}p{(\columnwidth - 10\tabcolsep) * \real{0.1486}}
  >{\raggedleft\arraybackslash}p{(\columnwidth - 10\tabcolsep) * \real{0.1622}}
  >{\raggedleft\arraybackslash}p{(\columnwidth - 10\tabcolsep) * \real{0.1757}}
  >{\raggedleft\arraybackslash}p{(\columnwidth - 10\tabcolsep) * \real{0.1351}}
  >{\raggedleft\arraybackslash}p{(\columnwidth - 10\tabcolsep) * \real{0.1757}}
  >{\raggedleft\arraybackslash}p{(\columnwidth - 10\tabcolsep) * \real{0.2027}}@{}}
\toprule\noalign{}
\begin{minipage}[b]{\linewidth}\raggedright
Group
\end{minipage} & \begin{minipage}[b]{\linewidth}\raggedleft
Estimate
\end{minipage} & \begin{minipage}[b]{\linewidth}\raggedleft
Scaled\_LS
\end{minipage} & \begin{minipage}[b]{\linewidth}\raggedleft
Day\_JK
\end{minipage} & \begin{minipage}[b]{\linewidth}\raggedleft
Period\_JK
\end{minipage} & \begin{minipage}[b]{\linewidth}\raggedleft
Reported\_SE
\end{minipage} \\
\midrule\noalign{}
\endhead
\bottomrule\noalign{}
\endlastfoot
Lunar & 1.0004 & 0.132 & 0.1218 & 0.147 & 0.132 \\
Near-12 h & 0.9251 & 0.2993 & 0.3649 & 0.3531 & 0.3649 \\
Diurnal & 0.5792 & 0.3887 & 0.5381 & 0.7735 & 0.5381 \\
\end{longtable}

\textbf{Extended Data Table 2 \textbar{} Conditional 95\% interval
coverage.} Each scenario has 1,000 simulated realizations per design,
with fixed observed sampling and weights. Values are percentages for the
larger-of-LS-and-day-jackknife rule, using multiplier 1.96. They assess
the specified scenarios, not unknown systematic effects.

\begin{longtable}[]{@{}llrrr@{}}
\toprule\noalign{}
scenario & design & Diurnal & Lunar & Near12h \\
\midrule\noalign{}
\endhead
\bottomrule\noalign{}
\endlastfoot
50\% OU 24 h, period & P8 & 97.9 & 98 & 97.3 \\
50\% OU 24 h, period & S54 & 96.4 & 95.1 & 95.4 \\
50\% OU 6 h, segment & P8 & 94.7 & 95.7 & 96 \\
50\% OU 6 h, segment & S54 & 95.8 & 95.6 & 96.3 \\
Student t (df 5) & P8 & 96.8 & 96.3 & 96.2 \\
Student t (df 5) & S54 & 96.2 & 96.1 & 96.8 \\
Gaussian & P8 & 95.7 & 96.2 & 95.9 \\
Gaussian & S54 & 95.5 & 97.4 & 96.7 \\
\end{longtable}

\textbf{Extended Data Table 3 \textbar{} Astronomical constituent
guide.} Names specify forcing and frequency; one response coefficient
multiplies each complete group with its internal ratios and phases
fixed.

\begin{longtable}[]{@{}lrll@{}}
\toprule\noalign{}
Symbol & Period\_h & Meaning & Group \\
\midrule\noalign{}
\endhead
\bottomrule\noalign{}
\endlastfoot
M2 & 12.421 & Principal lunar semidiurnal & Lunar \\
N2 & 12.658 & Lunar elliptic semidiurnal & Lunar \\
S2 & 12 & Principal solar semidiurnal & Near-12 h \\
K2 & 11.967 & Lunisolar semidiurnal & Near-12 h \\
K1 & 23.934 & Lunisolar diurnal & Diurnal \\
O1 & 25.819 & Principal lunar diurnal & Diurnal \\
P1 & 24.066 & Principal solar diurnal & Diurnal \\
Q1 & 26.868 & Lunar elliptic diurnal & Diurnal \\
\end{longtable}

\end{document}
